\scanpage{7}
% L1
On a un peu salopé la présentation, \deleted{\unsure{p}} qui suppose
% L2
pour le moins \deleted{\unsure{la}} car. \unsure{rés.} $\ne2$ (sans parler d'une
% L3
``théorie des angles'' sur schéma quelconque, que je
% L4
n'ai pas bien explicitée. Mais il est clair par
% L5
principe du prolongement que si la formule \inserted{sur l'invariant} est
% L6
prouvée en \deleted{car $2$.} $0$ (où on se sent sur terrain
% L7
solide) elle l'est en car. quelconque\ldots

% L8
Comme on est dans le cas de forme \inserted{quad} (\unsure{invariante non|invariante à hom}
% L9
id.\ nulle sur \deleted{\unsure{la}}\inserted{aucune} fibre) lisse, on doit bien avoir
% L10
$(\alpha_0+2)(\alpha_0-2)(\alpha_0+\alpha_1)$ inv. Inversement, si on part
% L11
d'un couple $(\alpha_0,\alpha_1)$ satisfaisant à ces conditions,
% L12
on sait lui associer un polyèdre régulier,
% L13
avec $\unsure{f_V|f_Y|\text{f.q.}}$ lisse, d'où représentation de $G_2$, correspondant
% L14
à ces valeurs. \deleted{On sait que la cat.}
% L15
\deleted{A) est}

% L16
\uline{Corollaire} \deleted{\unsure{Soient}}\unsure{Soit|Soient} $\unsure{\nu_0,\nu_1|v_0,v_1}$ deux entiers $\geq3$ \inserted{ou infinis}, et supposons
% L17
que l'on ait dans $\Gamma$ un syst\textsuperscript{e} complet de
% L18
relations
% L19
\[\tau_0^2=\tau_1^2=\tau_2^2=(\tau_0\tau_2)^2=(\tau_0\tau_1)^{\nu_0}=(\tau_{\unsure{1|0}}\tau_2)^{\nu_1}=1\]
% L20
i.e. $\Gamma=\Gamma_{\nu_0,\nu_1}$.
% slanted
% L20a
\margin{Supposons pour $i=0,1$, que $\nu_i\cdot1_S$ soit inv. sur $S$ \uline{ou} $\nu_i$ premier (impair puisque $\nu_i\geq3$)}
% L21
\deleted{Alors l'image par le fonc. des facteurs \unsure{de $C'$}} Soit \deleted{\unsure{le $G$ la sous}} $C(\nu_0,\nu_1)$
% L22
la sous-catégorie de $C$ formée des
% L23
$(G,\varphi)$ \inserted{($G$ forme de $\underline{GP}(1)_S$, $\varphi:\Gamma\longrightarrow G(S)$ hom)} tels que l'on ait
% L24
\deleted{$\unsure{\mathrm{i})}\quad\rho_0^{\nu_0}=\rho_1^{\nu_1}=1,$ \unsure{i.e., $\rho$}.}

% L25
a) $\forall s\in S$, $\rho_{0s}$ est d'ordre $\nu_0$, $\rho_{1s}$ d'ordre $\nu_1$ \uline{\unsure{exact\textsuperscript{ement}|exactem\textsuperscript{t}|exactement}}

% L26
b) $\forall s\in S$, \inserted{$\rho_0$ et $\rho_1$ ne commutent pas. \unsure{(en principe)}} \deleted{$\Gamma_s$ n'est pas isom. à un $\mathbb D_n$ \unsure{ni|ni à} $\unsure{\mathbb D_\infty}$}
% L27-30
\deleted{$\mathbb D_\infty\times\mathbb Z_{/2}$ [\uline{NB} Il revient au m à d'exiger que $\rho_{0s}$ et $\rho_{1s}$, ou $\rho_{0s}$ et $\sigma_s$, ne commutent pas, \deleted{et} il \uline{suffit} pour cela, compte tenu de a), que l'on ait : on n'a pas $\nu_0$ impair et $\nu_1=2\nu_0$ ou l'inverse]}
% archive sheet number: 6
